\documentclass[12pt,a4paper]{article}
\usepackage{newtxtext,newtxmath}
\usepackage[utf8]{inputenc}
\usepackage[T1]{fontenc}
\usepackage[margin=2.5cm]{geometry}
\usepackage[
backend=biber,
style=apa
]{biblatex}

\addbibresource{referencias.bib}
\usepackage{xcolor}
\usepackage{setspace}
\usepackage{csquotes}
\usepackage{titlesec}

% ---------------------------------------------------------
% DEFINICIÓN DE COLORES
% ---------------------------------------------------------
\definecolor{azulnoche}{RGB}{0,40,90}
\definecolor{celesteclaro}{RGB}{224,240,255}
\definecolor{grisclaro}{RGB}{245,245,245}
\definecolor{darkgreen}{rgb}{0.533, 0.486, 0.184}
\definecolor{golden}{RGB}{197,160,89}
\definecolor{goldford}{RGB}{143,113,54}
\definecolor{rojo}{RGB}{166,25,46}

% ---------------------------------------------------------
% CONFIGURACIÓN DE TÍTULOS
% ---------------------------------------------------------
\titleformat{\section}
{\Huge\bfseries\color{azulnoche}}
{\thesection}
{1em}
{}

\titleformat{\subsection}
{\Large\bfseries\color{azulnoche}}
{\thesubsection}
{1em}
{}

\titleformat{\subsubsection}
{\large\bfseries\color{azulnoche}}
{\thesubsubsection}
{1em}
{}

\usepackage{colortbl}
\usepackage{graphicx}
\usepackage{caption}
\captionsetup{compatibility=false}
\usepackage{booktabs}
\usepackage{tocloft}
\usepackage[spanish,es-tabla]{babel}
\usepackage{hyperref}
\usepackage{float}
\usepackage{tabularx}
\usepackage{array}
\usepackage{listings}

% ---------------------------------------------------------
% CONFIGURACIÓN DE LISTINGS (CÓDIGO JAVA)
% ---------------------------------------------------------
\lstset{
    language=Java,
    basicstyle=\ttfamily\small,
    keywordstyle=\color{azulnoche}\bfseries,
    commentstyle=\color{darkgreen}\itshape,
    stringstyle=\color{rojo},
    numbers=left,
    numberstyle=\tiny\color{gray},
    stepnumber=1,
    numbersep=8pt,
    frame=single,
    rulecolor=\color{azulnoche},
    backgroundcolor=\color{grisclaro},
    breaklines=true,
    showstringspaces=false,
    tabsize=4
}

\hypersetup{
    colorlinks=true,
    linkcolor=azulnoche,
    urlcolor=azulnoche,
    citecolor=azulnoche
}

\setlength{\parindent}{1.5em}
\onehalfspacing
\setlength{\cftbeforesecskip}{6pt}
\renewcommand{\contentsname}{Índice General}

% ---------------------------------------------------------
% ESTILOS DEL ÍNDICE CON COLORES RESPETADOS
% ---------------------------------------------------------
\renewcommand{\cftsecfont}{\color{azulnoche}\bfseries}
\renewcommand{\cftsubsecfont}{\color{azulnoche}}
\renewcommand{\cftsubsubsecfont}{\color{azulnoche}\itshape}

\renewcommand{\cftsecpagefont}{\color{azulnoche}\bfseries}
\renewcommand{\cftsubsecpagefont}{\color{azulnoche}}
\renewcommand{\cftsubsubsecpagefont}{\color{azulnoche}}

\renewcommand{\cftsecpresnum}{\color{azulnoche}}
\renewcommand{\cftsubsecpresnum}{\color{azulnoche}}
\renewcommand{\cftsubsubsecpresnum}{\color{azulnoche}}

\renewcommand{\cftsecaftersnum}{\color{azulnoche}}
\renewcommand{\cftsubsecaftersnum}{\color{azulnoche}}
\renewcommand{\cftsubsubsecaftersnum}{\color{azulnoche}}

\begin{document}

% =========================================================
% CARÁTULA
% =========================================================

\begin{titlepage}
\centering
{\color{golden}\bfseries\large UNIVERSIDAD NACIONAL MAYOR DE SAN MARCOS}\\
\vspace{0.2cm}
{\color{golden}\large\textbf{Universidad del Perú, Decana de América}}\\
\vspace{0.6cm}
\includegraphics[width=0.5\textwidth]{LOGOUNMSM.png}\\
\vspace{0.5cm}
{\color{goldford}\scshape\Large Facultad de Ingeniería de Sistemas e Informática\\}
{\color{goldford}\scshape\large Escuela Profesional de Ingeniería de Sistemas\\}
\vspace{0.6cm}
\color{rojo}\hrule height 0.3mm\vspace{0.5cm}
{\color{rojo}\Large\bfseries Semana 2 - POO\vspace{0.56cm}}\\
\color{rojo}\hrule height 0.3mm\vspace{0.6cm}
{\color{rojo}\Large\textbf{Curso:}}\\
\vspace{0.2cm}
{\color{goldford}\Large Programación de Computadoras II}\\
\vspace{0.7cm}
{\color{rojo}\Large\textbf{Docente:}}\\
\vspace{0.2cm}
{\color{goldford}\Large Javier Elmer Cabrera Díaz}\\
\vspace{0.7cm}
{\color{rojo}\Large\textbf{Autor:}}\\
\vspace{0.2cm}
{\color{goldford}\Large Mamani Huahualuque, Airton Abedal}\\
\vspace{1.5cm}
{\color{goldford}\Large\textbf{Lima, 8 de Septiembre del 2026}}\\
\end{titlepage}

% =========================================================
% ÍNDICE
% =========================================================

\newpage
\pagenumbering{arabic}
\setcounter{page}{1}

\tableofcontents

\newpage

% =========================================================
% DOCUMENTO
% =========================================================

\section{Introducción}

El paradigma de la Programación Orientada a Objetos (POO) constituye la columna vertebral del desarrollo de software moderno, permitiendo abstraer entidades complejas del mundo real hacia estructuras de código modulares, mantenibles y reutilizables. El presente informe aborda el desarrollo práctico de dos tareas fundamentales en el aprendizaje del lenguaje Java: la implementación del principio de encapsulamiento mediante una clase geométrica (Tarea 2.1) y el modelado de diversos niveles de asociación entre objetos (Tarea 2.2).

El objetivo principal de esta entrega es demostrar el dominio técnico en la construcción de clases robustas que protejan su estado interno mediante modificadores de acceso, métodos accesores/modificadores y validaciones de dominio. Asimismo, se ilustran de forma gráfica y conceptual los tres tipos de relaciones estructurales entre clases: \textbf{Asociación Simple}, \textbf{Agregación} y \textbf{Composición}, destacando las diferencias de acoplamiento y ciclo de vida de los componentes en cada caso.

\section{Desarrollo y Código Fuente}

\subsection{Tarea 2.1: Clase Rectangulo y Encapsulamiento}

Para la Tarea 2.1, se diseñó la clase \texttt{Rectangulo} aplicando encapsulamiento estricto. Los atributos \texttt{ancho} y \texttt{alto} se definieron con visibilidad privada, forzando a que cualquier modificación se realice a través de sus respectivos métodos \emph{setter}, los cuales garantizan valores mayores a cero. Se implementaron dos constructores (por defecto y parametrizado) y métodos para el cálculo de área y perímetro.

\begin{figure}[H]
    \centering
    \includegraphics[width=0.85\textwidth]{java1.jpg}
    \caption{Código fuente de la clase Rectangulo.java.}
    \label{fig:rectangulo_java}
    \vspace{0.2cm}
    \small{\textbf{Nota:} Implementación de atributos privados con validación en setters y sobrescritura de toString().}
\end{figure}

\begin{figure}[H]
    \centering
    \includegraphics[width=0.85\textwidth]{java2.jpg}
    \caption{Clase de prueba MainRectangulo.java.}
    \label{fig:main_rectangulo}
    \vspace{0.2cm}
    \small{\textbf{Nota:} Métodos de prueba para verificar la creación de objetos por defecto, parametrizados y modificación de atributos.}
\end{figure}

\subsection{Tarea 2.2: Relaciones de Asociación, Agregación y Composición}

En la Tarea 2.2 se representan las tres variantes de interacción entre objetos en Java:

\begin{itemize}
    \item \textbf{Asociación Simple (Vuelo - Aeropuerto):} Los aeropuertos existen de forma totalmente independiente del vuelo. Un vuelo referencia un aeropuerto de origen y uno de destino sin ser dueño de su ciclo de vida.
    \item \textbf{Agregación (Almacén - Montacargas):} Relación 'tiene-un' con acoplamiento débil. Un almacén utiliza un montacargas, pero si el almacén se destruye, el montacargas puede reasignarse a otro almacén.
    \item \textbf{Composición (Póliza - Cobertura):} Relación 'parte-de' con acoplamiento fuerte. La cobertura se instancia internamente dentro de la póliza de seguro. Si la póliza deja de existir, su cobertura asociada también se destruye.
\end{itemize}

\begin{figure}[H]
    \centering
    \includegraphics[width=0.85\textwidth]{java3.jpg}
    \caption{Clase Aeropuerto.java (Asociación Simple).}
    \label{fig:aeropuerto}
    \vspace{0.2cm}
    \small{\textbf{Nota:} Definición de la entidad independiente Aeropuerto con atributos de nombre y ciudad.}
\end{figure}

\begin{figure}[H]
    \centering
    \includegraphics[width=0.85\textwidth]{java4.jpg}
    \caption{Clase Vuelo.java (Asociación Simple).}
    \label{fig:vuelo}
    \vspace{0.2cm}
    \small{\textbf{Nota:} Incorporación de atributos de tipo Aeropuerto para relacionar origen y destino.}
\end{figure}

\begin{figure}[H]
    \centering
    \includegraphics[width=0.85\textwidth]{java5.jpg}
    \caption{Clase Montacargas.java (Agregación).}
    \label{fig:montacargas}
    \vspace{0.2cm}
    \small{\textbf{Nota:} Representación del recurso externo Montacargas que puede operar de forma independiente.}
\end{figure}

\begin{figure}[H]
    \centering
    \includegraphics[width=0.85\textwidth]{java6.jpg}
    \caption{Clase Almacen.java (Agregación).}
    \label{fig:almacen}
    \vspace{0.2cm}
    \small{\textbf{Nota:} Relación de agregación mediante la asignación externa de un objeto Montacargas.}
\end{figure}

\begin{figure}[H]
    \centering
    \includegraphics[width=0.85\textwidth]{java7.jpg}
    \caption{Clase Cobertura.java (Composición).}
    \label{fig:cobertura}
    \vspace{0.2cm}
    \small{\textbf{Nota:} Clase dependiente que representa las características de una cobertura de seguro.}
\end{figure}

\begin{figure}[H]
    \centering
    \includegraphics[width=0.85\textwidth]{java8.jpg}
    \caption{Clase PolizaSeguro.java (Composición).}
    \label{fig:poliza}
    \vspace{0.2cm}
    \small{\textbf{Nota:} Instanciación directa e interna del objeto Cobertura dentro del constructor de PolizaSeguro.}
\end{figure}

\begin{figure}[H]
    \centering
    \includegraphics[width=0.85\textwidth]{java9.jpg}
    \caption{Ejecución principal MainAsociaciones.java.}
    \label{fig:main_asociaciones}
    \vspace{0.2cm}
    \small{\textbf{Nota:} Programa principal que ejecuta las tres relaciones entre objetos y muestra el resultado en la consola.}
\end{figure}

\newpage

\section{Conclusiones}

\begin{itemize}
    \item \textbf{Garantía de Integridad mediante Encapsulamiento:} La restricción del acceso directo a los atributos de un objeto a través de modificadores \texttt{private} e interfaces de acceso (\emph{getters} y \emph{setters}) asegura que los datos no puedan tomar estados inválidos o inconsistentes dentro del sistema.
    \item \textbf{Flexibilidad y Reutilización de Código:} Al comprender la diferencia entre Agregación y Composición, el desarrollador puede tomar decisiones de arquitectura adecuadas: optar por agregación cuando se requiera compartir objetos entre distintos módulos, y composición cuando la existencia de un objeto secundario dependa indefectiblemente del objeto contenedor.
    \item \textbf{Mantenibilidad y Diseño Limpio:} El uso de la anotación \texttt{@Override} para rediseñar métodos estándar como \texttt{toString()} simplifica significativamente el proceso de depuración y salida de datos, promoviendo un código legible, autocontenido y estándar bajo las mejores prácticas de la industria Java.
\end{itemize}

\end{document}